% AVALIAÇÕES — Volume 2 (gerado por expandir_volume.py)

\chapter{Avaliações Diagnósticas}

\textbf{Sobre as avaliações:} no 2º ano, aplicam-se três avaliações: diagnóstica (início do ano), formativa (meio) e somativa (fim). Cada avaliação tem itens no formato das matrizes externas (CAEd/SPAECE) e gabarito comentado. As avaliações medem o desenvolvimento, não rotulam o aluno.

\section{Avaliação Diagnóstica — Início do 2º Ano}

\begin{avaliacao}
Objetivo: verificar a consolidação da base alfabética (sílabas simples, vogais) e a prontidão para os dígrafos. Aplicação individual ou em pequenos grupos, 30 minutos.
\end{avaliacao}

\noindent\textbf{Nome:} \underline{\hspace{4cm}} \quad \textbf{Data:} \underline{\hspace{2cm}}

\subsection*{Questão 1}
Leia em voz alta as sílabas e escreva ao lado o que você leu: MA  SA  TA  LA  PA.

\begin{center}
  \underline{\hspace{12cm}}
\end{center}

\subsection*{Questão 2}
Complete as palavras com a vogal que falta: m\_te (maté), s\_p\_ (sapo), f\_ra (fada).

\begin{center}
  \underline{\hspace{12cm}}
\end{center}

\subsection*{Questão 3}
Separe em sílabas e escreva o número de sílabas: BOLA, CASACA, MONTANHA.

\begin{center}
  \underline{\hspace{12cm}}
\end{center}

\subsection*{Questão 4}
Circule as palavras que terminam com o mesmo som: CHAVE e CAVE? CHUVA e BUVA? LHE e LE?

\begin{center}
  \underline{\hspace{12cm}}
\end{center}

\subsection*{Questão 5}
Escreva uma frase usando a palavra NINHO.

\begin{center}
  \underline{\hspace{12cm}}
\end{center}

\subsection*{Questão 6}
Leia o texto curto: 'O pato nada no lago.' e desenhe o que você entendeu.

\begin{center}
  \underline{\hspace{12cm}}
\end{center}

\clearpage

\section{Avaliação Formativa — Meio do 2º Ano}

\begin{avaliacao}
Objetivo: verificar os dígrafos CH, LH, NH e os encontros consonantais. 40 minutos, aplicação coletiva com leitura oral individual.
\end{avaliacao}

\noindent\textbf{Nome:} \underline{\hspace{4cm}} \quad \textbf{Data:} \underline{\hspace{2cm}}

\subsection*{Questão 1}
Forme sílabas com CH e escreva 3 palavras: CH+\_, CH+\_, CH+\_.

\begin{center}
  \underline{\hspace{12cm}}
\end{center}

\subsection*{Questão 2}
Complete com LH ou NH: coel\_o, ni\_ho, ora\_ha, gali\_a.

\begin{center}
  \underline{\hspace{12cm}}
\end{center}

\subsection*{Questão 3}
Circule os encontros consonantais e leia: BRA, CRA, DRA, FRA, GRA, PRA, TRA.

\begin{center}
  \underline{\hspace{12cm}}
\end{center}

\subsection*{Questão 4}
Escreva 3 palavras com BR e 3 palavras com TR.

\begin{center}
  \underline{\hspace{12cm}}
\end{center}

\subsection*{Questão 5}
Leia a frase e sublinhe as palavras com encontro consonantal: 'O prato tem fruta e o trem chega cedo.'

\begin{center}
  \underline{\hspace{12cm}}
\end{center}

\subsection*{Questão 6}
Separe em sílabas: PRATO, DRAGÃO, FLORESTA, GLOBO.

\begin{center}
  \underline{\hspace{12cm}}
\end{center}

\clearpage

\section{Avaliação Somativa — Fim do 2º Ano}

\begin{avaliacao}
Objetivo: verificar as grafias RR, SS, S/Z, Ç, AM/AN e X, além da fluência em texto curto. 50 minutos.
\end{avaliacao}

\noindent\textbf{Nome:} \underline{\hspace{4cm}} \quad \textbf{Data:} \underline{\hspace{2cm}}

\subsection*{Questão 1}
Complete com RR ou R: ca\_o, te\_a, mo\_o, co\_o (corro).

\begin{center}
  \underline{\hspace{12cm}}
\end{center}

\subsection*{Questão 2}
Complete com S, SS ou Z: ca\_a (casa), ca\_a (caça), pa\_ar (passar), me\_a (mesa).

\begin{center}
  \underline{\hspace{12cm}}
\end{center}

\subsection*{Questão 3}
Complete com Ç, C ou S: cora\_ão, \_apéu, \_alinha, \_açarola.

\begin{center}
  \underline{\hspace{12cm}}
\end{center}

\subsection*{Questão 4}
Complete com AM ou AN: cant\_\_, brinc\_\_, s\_\_gue, c\_\_po.

\begin{center}
  \underline{\hspace{12cm}}
\end{center}

\subsection*{Questão 5}
Leia o texto: 'O carro passou na estrada. A professora sorriu. O pássaro cantou.' e responda: o que passou na estrada?

\begin{center}
  \underline{\hspace{12cm}}
\end{center}

\subsection*{Questão 6}
Escreva uma história curta (4 frases) usando pelo menos uma palavra com CH, uma com NH e uma com RR.

\begin{center}
  \underline{\hspace{12cm}}
\end{center}

\clearpage

\section{Avaliação de Leitura e Compreensão — 3º Trimestre}

\begin{avaliacao}
Objetivo: verificar fluência leitora (precisão, ritmo, expressividade) e compreensão de texto curto. Aplicação individual da leitura em voz alta; as questões são respondidas por escrito.
\end{avaliacao}

\noindent\textbf{Nome:} \underline{\hspace{4cm}} \quad \textbf{Data:} \underline{\hspace{2cm}}

\subsection*{Questão 1}
Leia em voz alta o texto: 'O carro do vizinho parou na chuva. A professora abriu a porta e sorriu. O menino tirou o chapéu e entrou na escola.'

\begin{center}
  \underline{\hspace{12cm}}
\end{center}

\subsection*{Questão 2}
Quem chegou à escola?

\begin{center}
  \underline{\hspace{12cm}}
\end{center}

\subsection*{Questão 3}
Onde o carro parou?

\begin{center}
  \underline{\hspace{12cm}}
\end{center}

\subsection*{Questão 4}
O que a professora fez ao abrir a porta?

\begin{center}
  \underline{\hspace{12cm}}
\end{center}

\subsection*{Questão 5}
Retire do texto uma palavra com CH.

\begin{center}
  \underline{\hspace{12cm}}
\end{center}

\subsection*{Questão 6}
Retire do texto uma palavra com RR.

\begin{center}
  \underline{\hspace{12cm}}
\end{center}

\subsection*{Questão 7}
Quantas palavras tem a primeira frase do texto?

\begin{center}
  \underline{\hspace{12cm}}
\end{center}

\subsection*{Questão 8}
Qual é a melhor nova frase para continuar a história? Escreva uma.

\begin{center}
  \underline{\hspace{12cm}}
\end{center}

\clearpage

\section{Avaliação de Produção Escrita — Final do Ano}

\begin{avaliacao}
Objetivo: verificar a produção de texto com estrutura (início, meio, fim), ortografia das grafias do ano e legibilidade. 45 minutos; o rascunho é corrigido antes da versão final.
\end{avaliacao}

\noindent\textbf{Nome:} \underline{\hspace{4cm}} \quad \textbf{Data:} \underline{\hspace{2cm}}

\subsection*{Questão 1}
Planeje: escolha o tema do seu texto (um dia de chuva, a visita à biblioteca, o passeio com o coelho) e escreva 3 palavras que vai usar.

\begin{center}
  \underline{\hspace{12cm}}
\end{center}

\subsection*{Questão 2}
Escreva o rascunho do texto com, no mínimo, 4 frases. Use uma palavra com CH, uma com NH e uma com RR.

\begin{center}
  \underline{\hspace{12cm}}
\end{center}

\subsection*{Questão 3}
Releia o rascunho: circule as palavras com CH, LH, NH, RR, SS, S entre vogais, Ç, AM/AN e X que você usou.

\begin{center}
  \underline{\hspace{12cm}}
\end{center}

\subsection*{Questão 4}
Corrija a ortografia: troque as grafias erradas usando a regra que aprendeu no ano.

\begin{center}
  \underline{\hspace{12cm}}
\end{center}

\subsection*{Questão 5}
Escreva a versão final em letra cursiva, com margem e parágrafo.

\begin{center}
  \underline{\hspace{12cm}}
\end{center}

\subsection*{Questão 6}
Verifique: comecei com letra maiúscula? Terminei com ponto final? Todo parágrafo tem seu espaço?

\begin{center}
  \underline{\hspace{12cm}}
\end{center}

\subsection*{Questão 7}
Ilustre o seu texto com um desenho.

\begin{center}
  \underline{\hspace{12cm}}
\end{center}

\subsection*{Questão 8}
Apresente o texto para a turma em voz alta, com fluência e expressividade.

\begin{center}
  \underline{\hspace{12cm}}
\end{center}

\clearpage

\section{Avaliação de Ortografia — 2º Semestre}

\begin{avaliacao}
Objetivo: verificar o domínio das grafias do 2º semestre (AM/AN, X, S/Z/Ç consolidados) e a justificativa das regras. Aplicação coletiva, 40 minutos.
\end{avaliacao}

\noindent\textbf{Nome:} \underline{\hspace{4cm}} \quad \textbf{Data:} \underline{\hspace{2cm}}

\subsection*{Questão 1}
Complete com AM ou AN: eles cant\_\_, o s\_\_gue, o c\_\_po, o t\_\_bém? (também).

\begin{center}
  \underline{\hspace{12cm}}
\end{center}

\subsection*{Questão 2}
Complete com X, CH ou SS: pei\_e, \_uva, pa\_ar, \_adrez.

\begin{center}
  \underline{\hspace{12cm}}
\end{center}

\subsection*{Questão 3}
Complete com S ou Z: ca\_a, \_ebra, me\_a, ra\_ão.

\begin{center}
  \underline{\hspace{12cm}}
\end{center}

\subsection*{Questão 4}
Complete com Ç ou SS: cora\_ão, po\_o, a\_úcar, \_apéu.

\begin{center}
  \underline{\hspace{12cm}}
\end{center}

\subsection*{Questão 5}
Explique com suas palavras: por que CASA tem som de Z?

\begin{center}
  \underline{\hspace{12cm}}
\end{center}

\subsection*{Questão 6}
Escreva 3 palavras com X e diga o som de cada uma.

\begin{center}
  \underline{\hspace{12cm}}
\end{center}

\subsection*{Questão 7}
Escreva o plural de: CANÇÃO (CANÇÕES), LIÇÃO (LIÇÕES), PÃO (PÃES).

\begin{center}
  \underline{\hspace{12cm}}
\end{center}

\subsection*{Questão 8}
Corrija a frase: 'O caro paso na rua suja' (O carro passou na rua suja).

\begin{center}
  \underline{\hspace{12cm}}
\end{center}

\clearpage

\section{Avaliação de Fluência Leitora — Final do Ano}

\begin{avaliacao}
Objetivo: medir a fluência (precisão, ritmo, expressividade) em texto conhecido e desconhecido. Aplicação individual: leitura de 1 minuto com registro de palavras lidas corretamente.
\end{avaliacao}

\noindent\textbf{Nome:} \underline{\hspace{4cm}} \quad \textbf{Data:} \underline{\hspace{2cm}}

\subsection*{Questão 1}
Leia o texto 1 (conhecido) em voz alta por 1 minuto: 'O coelho fugiu do jardim. A chuva começou a cair. Ele correu para o ninho e se escondeu até o sol voltar.'

\begin{center}
  \underline{\hspace{12cm}}
\end{center}

\subsection*{Questão 2}
Leia o texto 2 (desconhecido) em voz alta por 1 minuto: 'A biblioteca da escola ganhou uma estante nova. As crianças escolheram histórias de fadas e de dragões. Todos queriam ler primeiro o livro do coelho aventureiro.'

\begin{center}
  \underline{\hspace{12cm}}
\end{center}

\subsection*{Questão 3}
O professor registra: palavras lidas corretamente no texto 1 e no texto 2.

\begin{center}
  \underline{\hspace{12cm}}
\end{center}

\subsection*{Questão 4}
Responda oralmente: de quem fala o texto 1? O que aconteceu depois da chuva?

\begin{center}
  \underline{\hspace{12cm}}
\end{center}

\subsection*{Questão 5}
Responda oralmente: o que a biblioteca ganhou? Qual livro todos queriam ler?

\begin{center}
  \underline{\hspace{12cm}}
\end{center}

\subsection*{Questão 6}
Leia com expressividade a frase: 'Todos queriam ler primeiro o livro do coelho aventureiro!'

\begin{center}
  \underline{\hspace{12cm}}
\end{center}

\subsection*{Questão 7}
Autoavaliação: minha leitura ficou mais rápida e bonita do que no início do ano?

\begin{center}
  \underline{\hspace{12cm}}
\end{center}

\clearpage

\section{Avaliação Integrada Final do Ano (2º Ano)}

\begin{avaliacao}
Objetivo: verificar a consolidação das metas do ano (leitura fluente, ortografia das grafias trabalhadas, produção de gêneros e compreensão). Aplicação em duas sessões de 40 minutos, com pausa.
\end{avaliacao}

\noindent\textbf{Nome:} \underline{\hspace{4cm}} \quad \textbf{Data:} \underline{\hspace{2cm}}

\subsection*{Questão 1}
Leia o texto: 'O menino soltava pipas no campo. O vento levou a pipa para longe. O menino não ficou triste: correu atrás dela e a encontrou perto da cerca.' Quantas frases tem o texto?

\begin{center}
  \underline{\hspace{12cm}}
\end{center}

\subsection*{Questão 2}
Quem soltava pipas? Para onde o vento levou a pipa? Onde o menino encontrou a pipa?

\begin{center}
  \underline{\hspace{12cm}}
\end{center}

\subsection*{Questão 3}
Retire do texto: uma palavra com TR, uma palavra com ENCONTRO DOIS? (campo tem AN).

\begin{center}
  \underline{\hspace{12cm}}
\end{center}

\subsection*{Questão 4}
Escreva o final da história com 2 frases usando uma palavra com CH e uma com NH.

\begin{center}
  \underline{\hspace{12cm}}
\end{center}

\subsection*{Questão 5}
Complete: a pipa \_ooa no céu (voa). O menino \_orreu (correu). A \_erca é de madeira (cerca).

\begin{center}
  \underline{\hspace{12cm}}
\end{center}

\subsection*{Questão 6}
Escreva um bilhete para um amigo combinando de soltar pipas no sábado.

\begin{center}
  \underline{\hspace{12cm}}
\end{center}

\subsection*{Questão 7}
Separe em sílabas: PIPAS, CAMPO, CERCA, VENTO.

\begin{center}
  \underline{\hspace{12cm}}
\end{center}

\subsection*{Questão 8}
Autoavaliação: o que eu aprendi de melhor neste ano? O que ainda quero melhorar?

\begin{center}
  \underline{\hspace{12cm}}
\end{center}

\clearpage

\chapter{Gabaritos e Matriz de Correção}

\section{Gabarito --- Avaliação Diagnóstica — Início do 2º Ano}

\begin{center}
\renewcommand{\arraystretch}{1.5}
\begin{tabular}{lll}
\toprule
**Item** & **Resposta** & **Justificativa** \\
\midrule

Q1 & MA, SA, TA, LA, PA & Decodificação de sílabas diretas consolidadas \\

Q2 & maté/sapo/fada & Identificação de vogais em palavras \\

Q3 & 2, 3, 3 sílabas & Segmentação silábica \\

Q4 & CHAVE (som /ch/) & Consciência do som-alvo \\

Q5 & Frase com NINHO & Produção escrita com palavras do ano \\

Q6 & Desenho coerente & Compreensão leitora inicial \\

\bottomrule
\end{tabular}
\end{center}

\section{Gabarito --- Avaliação Formativa — Meio do 2º Ano}

\begin{center}
\renewcommand{\arraystretch}{1.5}
\begin{tabular}{lll}
\toprule
**Item** & **Resposta** & **Justificativa** \\
\midrule

Q1 & CHA CHE CHI CHO CHU + palavras & Produção com dígrafo CH \\

Q2 & coelho, ninho, orelha, galinha & Ortografia LH/NH \\

Q3 & Encontros identificados & Reconhecimento de encontros \\

Q4 & Palavras com BR e TR & Produção de palavras com encontros \\

Q5 & PRATO, FRUTA, TREM & Leitura e identificação em contexto \\

Q6 & PRA-TO, DRA-GÃO, FLO-RES-TA, GLO-BO & Segmentação de encontros \\

\bottomrule
\end{tabular}
\end{center}

\section{Gabarito --- Avaliação Somativa — Fim do 2º Ano}

\begin{center}
\renewcommand{\arraystretch}{1.5}
\begin{tabular}{lll}
\toprule
**Item** & **Resposta** & **Justificativa** \\
\midrule

Q1 & carro, terra, morro, corro & Ortografia RR \\

Q2 & casa, caça, passar, mesa & S entre vogais, SS e Z \\

Q3 & coração, chapéu, galinha, caçarola & Ç e dígrafos \\

Q4 & cantam, brincam, sangue, campo & AM/AN \\

Q5 & O carro & Compreensão leitora \\

Q6 & Produção com grafias do ano & Escrita espontânea \\

\bottomrule
\end{tabular}
\end{center}

\section{Gabarito --- Avaliação de Leitura e Compreensão — 3º Trimestre}

\begin{center}
\renewcommand{\arraystretch}{1.5}
\begin{tabular}{lll}
\toprule
**Item** & **Resposta** & **Justificativa** \\
\midrule

Q1 & Leitura oral com precisão e ritmo & Fluência \\

Q2 & O menino & Localização de informação \\

Q3 & Na chuva & Localização de informação \\

Q4 & Sorriu & Compreensão \\

Q5 & chapéu ou chuva & Identificação de grafia em texto \\

Q6 & carro & Identificação de grafia em texto \\

Q7 & 7 palavras & Análise linguística \\

Q8 & Resposta coerente com o texto & Produção de continuação \\

\bottomrule
\end{tabular}
\end{center}

\section{Gabarito --- Avaliação de Produção Escrita — Final do Ano}

\begin{center}
\renewcommand{\arraystretch}{1.5}
\begin{tabular}{lll}
\toprule
**Item** & **Resposta** & **Justificativa** \\
\midrule

Q1 & Planejamento com 3 palavras & Planejamento de escrita \\

Q2 & Texto com 4+ frases e grafias solicitadas & Produção textual \\

Q3 & Identificação das grafias no próprio texto & Metacognição ortográfica \\

Q4 & Correção ortográfica com justificativa & Autorregulação \\

Q5 & Versão final em cursiva legível & Caligrafia e estrutura \\

Q6 & Autoavaliação respondida & Autorrevisão \\

Q7 & Ilustração coerente & Integração multimodal \\

Q8 & Apresentação oral expressiva & Fluência e expressividade \\

\bottomrule
\end{tabular}
\end{center}

\section{Gabarito --- Avaliação de Ortografia — 2º Semestre}

\begin{center}
\renewcommand{\arraystretch}{1.5}
\begin{tabular}{lll}
\toprule
**Item** & **Resposta** & **Justificativa** \\
\midrule

Q1 & cantam, sangue, campo, também & AM/AN \\

Q2 & peixe, chuva, passar, xadrez & X, CH, SS \\

Q3 & casa, zebra, mesa, razão & S entre vogais e Z \\

Q4 & coração, poço, açúcar, chapéu & Ç e digrafos \\

Q5 & Regra do S entre vogais com som de Z & Justificativa ortográfica \\

Q6 & xadrez (/ch/), exame (/z/), táxi (/ks/) & Polissemia do X \\

Q7 & canções, lições, pães & Ditongos nasais \\

Q8 & carro, passou & Autocorreção ortográfica \\

\bottomrule
\end{tabular}
\end{center}

\section{Gabarito --- Avaliação de Fluência Leitora — Final do Ano}

\begin{center}
\renewcommand{\arraystretch}{1.5}
\begin{tabular}{lll}
\toprule
**Item** & **Resposta** & **Justificativa** \\
\midrule

Q1–Q2 & Registro de palavras/minuto & Medida de fluência \\

Q3 & Comparação dos dois textos & Precisão e ritmo \\

Q4 & Do coelho; ele se escondeu no ninho & Compreensão do texto conhecido \\

Q5 & Uma estante nova; o livro do coelho aventureiro & Compreensão do texto novo \\

Q6 & Leitura expressiva com entonação & Prosódia \\

Q7 & Reflexão sobre o progresso & Metacognição \\

\bottomrule
\end{tabular}
\end{center}

\section{Gabarito --- Avaliação Integrada Final do Ano (2º Ano)}

\begin{center}
\renewcommand{\arraystretch}{1.5}
\begin{tabular}{lll}
\toprule
**Item** & **Resposta** & **Justificativa** \\
\midrule

Q1 & 3 frases & Estrutura do texto \\

Q2 & O menino; para longe; perto da cerca & Localização de informações \\

Q3 & soltava? (TR); campo (AN) & Identificação de grafias \\

Q4 & Final coerente com CH e NH & Produção textual \\

Q5 & voa, correu, cerca & Ortografia \\

Q6 & Bilhete com estrutura (para quem, recado, de quem) & Gênero textual \\

Q7 & PI-PAS, CAM-PO, CER-CA, VEN-TO & Segmentação silábica \\

Q8 & Reflexão escrita & Metacognição \\

\bottomrule
\end{tabular}
\end{center}
